\chapter{Code Generation}
\label{ch:codegen}

% Source: src/zgen.c

\section{Overview}

The code generation phase (\texttt{zgen.c}) walks the type-checked AST and emits LLVM IR as text, which is then compiled to a native binary by invoking \texttt{clang} (or written to a \texttt{.ll} file when \texttt{--emit-llvm} is passed).

The main entry point is:

\begin{verbatim}
void generate(ZState *state, ZNode **nodes);
\end{verbatim}

\section{IR Emission Strategy}

Zinc uses the LLVM C API to build the module in memory, then serialises it with \texttt{LLVMPrintModuleToString} / \texttt{LLVMPrintModuleToFile}.

Key internal helpers:

\begin{description}[leftmargin=4em]
  \item[\texttt{genFunc}]         Emit a function definition, allocate stack slots for parameters.
  \item[\texttt{genFuncVars}]     Pre-scan the function body for variable declarations and emit \texttt{alloca}s at the entry block (required by LLVM's \texttt{mem2reg} pass).
  \item[\texttt{genExpr}]         Recursively emit an expression and return its \texttt{LLVMValueRef}.
  \item[\texttt{genStructLitInto}] Emit a struct literal into a pre-allocated destination pointer.
  \item[\texttt{putDestructuredPatternInStack}] Emit loads/stores for tuple or struct destructuring.
\end{description}

\section{Stack Allocation Model}

All local variables are allocated with \texttt{alloca} at the top of the entry block.
Loads and stores are emitted on use and assignment respectively, and LLVM's \texttt{mem2reg} pass promotes them to SSA registers.

\texttt{genFuncVars} performs a pre-pass that visits every \texttt{NODE\_VAR\_DEF} in the function body and emits the \texttt{alloca} before any other instruction. This guarantees that variables are allocated in the dominator block.

\begin{note}
  \texttt{genFuncVars} currently has no case for \texttt{NODE\_ARRAY\_LIT} inside a \zn{return} statement, which means \zn{return [...]} fails with ``Missing stack value''. See B05 in the bug tracker.
\end{note}

\section{Struct Layout}

Structs are lowered to LLVM \texttt{\{T1, T2, \ldots\}} aggregate types.
Field access is compiled to a \texttt{getelementptr} (GEP) instruction followed by a \texttt{load} or \texttt{store}.

For generic structs, the type arguments are substituted during monomorphization (semantic pass) so the codegen phase always sees a concrete struct type.

\section{Function Calls}

Regular calls are emitted with \texttt{LLVMBuildCall2}.
Receiver-method calls are sugar: the receiver expression is evaluated and passed as the last argument (after the regular parameter list).

\subsection{Static Method Calls (\texttt{T::method})}

Calls of the form \zn{Point::print()} resolve to a regular function call with no implicit receiver argument.

\section{Control Flow}

\begin{description}[leftmargin=4em]
  \item[\texttt{if}/\texttt{else}] Three basic blocks: \texttt{then}, \texttt{else} (optional), \texttt{merge}.
  \item[C-style \zn{for}]         \texttt{init} block -> \texttt{cond} block -> \texttt{body} block -> \texttt{incr} block -> \texttt{cond}.
  \item[While-style \zn{for}]     \texttt{cond} block -> \texttt{body} block -> \texttt{cond}.
  \item[\zn{break}/\zn{continue}] Unconditional branches to the loop's \texttt{exit} or \texttt{cond} block, tracked on a stack.
  \item[\zn{goto}/labels]         Named basic blocks created on first reference; resolved after the function body is complete.
  \item[\zn{defer}]               Deferred expressions are collected in a list and emitted in reverse order before every \zn{return} in the current scope.
\end{description}

\section{Linking}

After IR emission, \texttt{zlink.cpp} drives the linker via the LLVM \texttt{lld} library (or falls back to the system linker).
The standard library (\texttt{lib/}) is linked unconditionally.